\scanpage{15}
\mspage{8}
% letterhead
\begin{center}
ACADEMIE D'AIX-MARSEILLE\\
\rule{4em}{0.4pt}\\
INSPECTION ACADEMIQUE\\
DE VAUCLUSE\\
\rule{1.6em}{0.4pt}
\end{center}
% ll. 1–8
épimorphismes. Alors \deleted{l'\unsure{application|applicat.}} pour $u:b\longrightarrow b'$ donné, l'application
\[
\begin{aligned}
 u_*\deleted{\unsure{*|=}}=\widetilde\pi(u):\quad\widetilde\pi(b)&\longrightarrow\widetilde\pi(b'),\qquad E\longrightarrow\pi(u)(E),\\
 \rotatebox{90}{$=$}&\hspace{3em}\rotatebox{90}{$=$}\\
 \mathcal P^*(\pi(b))&\longrightarrow\mathcal P^*(\pi(b'))
\end{aligned}
\]
admet un adjoint à droite \deleted{$\unsure{\pi|u}$} $u^*$, savoir $E'\longrightarrow\pi(u)^{-1}(E')$, qui transforme un $E'$ \uline{non vide} en \deleted{$\unsure{F|I}$}$E$ non vide. (\unsure{because|beca-use} $\pi(b)\longrightarrow\pi(b')$ surjectif)
% ll. 9–15
\uline{Lemme 1} Soit $C\longrightarrow B$ une \uline{fibration}, dans $\mathrm{Cat}$, et \deleted{\unsure{suppos|suppos.}} soit $I$ dans $\mathrm{Cat}$ telle que les $\varprojlim$ de type $I$ existent dans $B$, et dans les fibres $C_b$. Alors elles existent dans $C$, et $C\longrightarrow B$ commute aux $\varprojlim$ de type $I$.
% ll. 16–19
\uline{Corollaire} Soit $\pi:B\longrightarrow\mathrm{Ens}$ un foncteur, \deleted{transformant flèches de $B$ en \unsure{épis|épi}}, soit
\diagram{\begin{tikzcd}[ampersand replacement=\&]
 \widetilde\pi:B\arrow[r,"\pi"]\&\mathrm{Ens}\arrow[r,"{\mathcal P^{\unsure{*|\epsilon}}}"]\&\mathrm{Ord}\arrow[d,"\text{can.}"]\\
 \&\&\mathrm{Cat}
\end{tikzcd}}
le composé, \deleted{\unsure{et|de}} Alors $C$
% archive sheet number: 14
